\documentclass{article}
\usepackage{amsmath,amsfonts,amssymb}
\usepackage{graphicx}
\usepackage{booktabs}
\title{LIA Experiment Results}
\author{Automated Experiment Pipeline}
\date{August 21, 2025}
\begin{document}
\maketitle
\section{Introduction}
This document contains the results from the Lorentz-Invariant Auction (LIA) mechanism experiments. The experiments compare three auction mechanisms across three different network topologies.\n\n\section{Experimental Setup}
The experiments were conducted across three network topologies:
\begin{itemize}
\item \textbf{STARLINK-200}: Satellite constellation with 200 nodes (delays: 1-50 ms)
\item \textbf{INTERNET-100}: Internet backbone with 100 nodes (delays: 0.1-300 ms)
\item \textbf{DSN-30}: Deep Space Network with 30 nodes (delays: 60-2847 seconds)
\end{itemize}
Three auction mechanisms were compared:
\begin{itemize}
\item \textbf{LIA}: Lorentz-Invariant Auction (novel mechanism)
\item \textbf{VCG}: Vickrey-Clarke-Groves (baseline)
\item \textbf{HoldBack}: HoldBack mechanism (delay-aware baseline)
\end{itemize}
\section{Results}
\begin{table}[h!]
\centering
\caption{LIA Experiment Results Summary}
\begin{tabular}{|l|l|c|c|c|}
\hline
Mechanism & Topology & Efficiency & Revenue & Runtime (s) \\
\hline
HoldBack & DSN-30 & 0.876 $\pm$ 0.044 & 0.627 $\pm$ 0.071 & 0.0000 \\
HoldBack & INTERNET-100 & 0.874 $\pm$ 0.043 & 0.621 $\pm$ 0.073 & 0.0000 \\
HoldBack & STARLINK-200 & 0.874 $\pm$ 0.042 & 0.626 $\pm$ 0.072 & 0.0000 \\
LIA & DSN-30 & 0.917 $\pm$ 0.038 & 0.749 $\pm$ 0.060 & 0.0000 \\
LIA & INTERNET-100 & 0.915 $\pm$ 0.037 & 0.750 $\pm$ 0.057 & 0.0001 \\
LIA & STARLINK-200 & 0.917 $\pm$ 0.038 & 0.750 $\pm$ 0.058 & 0.0001 \\
VCG & DSN-30 & 1.000 $\pm$ 0.000 & 0.673 $\pm$ 0.073 & 0.0000 \\
VCG & INTERNET-100 & 1.000 $\pm$ 0.000 & 0.676 $\pm$ 0.074 & 0.0000 \\
VCG & STARLINK-200 & 1.000 $\pm$ 0.000 & 0.673 $\pm$ 0.073 & 0.0000 \\
\hline
\end{tabular}
\label{tab:results}
\end{table}

\section{Conclusion}
The results demonstrate the performance characteristics of the LIA mechanism compared to traditional auction mechanisms in heterogeneous network environments.
\end{document}
